Several caveats temper the interpretation of the belief measure. The elicitation was unincentivised and took place after the delegation decision, so reports may reflect hedging across scenarios or post-hoc rationalisation rather than the beliefs that actually drove the choice. As a corollary, any within-subject correlation between beliefs and delegation cannot identify direction: we cannot distinguish whether subjects delegated \textit{because} they held these specific beliefs about punishment, or hold these beliefs \textit{because} they delegated. Hypotheticality compounds the noise: only two of the four \textit{(delegation, outcome)} scenarios remained payoff-relevant after the delegation decision in the \textit{Punishment} condition, and all four were purely hypothetical in the \textit{No-Punishment} condition. We therefore restrict the belief--delegation analysis to the \textit{Punishment} condition. The motivation is structural rather than statistical: the punishment-anticipation channel can only operate where punishment is possible, so beliefs about a counterfactual that cannot materialise do not enter Player~A's decision rule in the \textit{No-Punishment} arm. Pooling the two conditions would therefore constitute a category error rather than a gain in power, and the exercise that follows should accordingly be read as suggestive evidence on whether the channel is alive among those exposed to it, not as a mediation analysis of the treatment-level delegation gap (Result~1).
